\section{Definitions and examples of symmetric polynomials}

\subsection{Setup}

\begin{convention}
\label{conv.sf.KN}
Fix a commutative ring $K$. Fix an $N\in\mathbb{N}$.
Throughout this chapter, we will keep $K$ and $N$ fixed.
Let $S_N$ denote the symmetric group, i.e., the group of all
permutations of $[N] := \{1,2,\ldots,N\}$.
\end{convention}

\subsection{The polynomial ring and the symmetric group action}

\begin{definition}
\label{def.sf.PS}
\lean{AlgebraicCombinatorics.SymmetricPolynomials.P, AlgebraicCombinatorics.SymmetricPolynomials.permAction, AlgebraicCombinatorics.SymmetricPolynomials.IsSymm, AlgebraicCombinatorics.SymmetricPolynomials.S}
\leantarget
\leanok
\textbf{(a)} Let $\mathcal{P}$ be the polynomial ring
$K[x_1, x_2, \ldots, x_N]$ in $N$ variables over $K$. This
is not just a ring; it is a commutative $K$-algebra. \medskip

\textbf{(b)} The symmetric group $S_N$ acts on the set $\mathcal{P}$
according to the formula
\[
\sigma \cdot f = f[x_{\sigma(1)}, x_{\sigma(2)}, \ldots, x_{\sigma(N)}]
\quad \text{for any } \sigma \in S_N \text{ and any } f \in \mathcal{P}.
\]
Here, $f[a_1, a_2, \ldots, a_N]$ means the result of
substituting $a_1, a_2, \ldots, a_N$ for the indeterminates $x_1, x_2, \ldots, x_N$
in a polynomial $f \in \mathcal{P}$.

Roughly speaking, the group $S_N$ is thus acting on $\mathcal{P}$ by
permuting variables: A permutation $\sigma \in S_N$ transforms a polynomial
$f$ by substituting $x_{\sigma(i)}$ for each $x_i$.

Note that this action of $S_N$ on $\mathcal{P}$ is a well-defined group
action (as we will see in Proposition~\ref{prop.sf.SN-acts} below). \medskip

\textbf{(c)} A polynomial $f \in \mathcal{P}$ is said to be \emph{symmetric} if
it satisfies
\[
\sigma \cdot f = f \quad \text{for all } \sigma \in S_N.
\]

\textbf{(d)} We let $\mathcal{S}$ be the set of all symmetric polynomials
$f \in \mathcal{P}$.
\end{definition}

\subsection{The action is well-defined}

\begin{proposition}
\label{prop.sf.SN-acts}
\lean{AlgebraicCombinatorics.SymmetricPolynomials.permAction_id, AlgebraicCombinatorics.SymmetricPolynomials.permAction_mul, AlgebraicCombinatorics.SymmetricPolynomials.permMulAction}
\leantarget
\leanok
The action of $S_N$ on $\mathcal{P}$ is a
well-defined group action. In other words, the following holds: \medskip

\textbf{(a)} We have $\operatorname{id}_{[N]} \cdot f = f$ for every $f \in \mathcal{P}$. \medskip

\textbf{(b)} We have $(\sigma\tau) \cdot f = \sigma \cdot (\tau \cdot f)$
for every $\sigma, \tau \in S_N$ and $f \in \mathcal{P}$.
\end{proposition}

\begin{proof}
\textbf{(a)} If $f \in \mathcal{P}$, then the definition of
$\operatorname{id}_{[N]} \cdot f$ yields
\[
\operatorname{id}_{[N]} \cdot f
= f[x_{\operatorname{id}(1)}, x_{\operatorname{id}(2)}, \ldots, x_{\operatorname{id}(N)}]
= f[x_1, x_2, \ldots, x_N]
= f.
\]
This proves part \textbf{(a)}. \medskip

\textbf{(b)} Let $\sigma, \tau \in S_N$ and $f \in \mathcal{P}$. The definition
of $\sigma \cdot (\tau \cdot f)$ yields
\begin{equation}
\sigma \cdot (\tau \cdot f)
= (\tau \cdot f)[x_{\sigma(1)}, x_{\sigma(2)}, \ldots, x_{\sigma(N)}].
\label{pf.prop.sf.SN-acts.b.1}
\end{equation}
Write the polynomial $f$ in the form
\begin{equation}
f = \sum_{(a_1, a_2, \ldots, a_N) \in \mathbb{N}^N} f_{(a_1, a_2, \ldots, a_N)} x_1^{a_1} x_2^{a_2} \cdots x_N^{a_N}.
\label{pf.prop.sf.SN-acts.b.f=}
\end{equation}
The definition of $\tau \cdot f$ yields
\[
\tau \cdot f = \sum_{(a_1, \ldots, a_N) \in \mathbb{N}^N} f_{(a_1, \ldots, a_N)} x_{\tau(1)}^{a_1} x_{\tau(2)}^{a_2} \cdots x_{\tau(N)}^{a_N}.
\]
Substituting $x_{\sigma(1)}, x_{\sigma(2)}, \ldots, x_{\sigma(N)}$
for $x_1, x_2, \ldots, x_N$ on both sides, we obtain
\begin{align*}
(\tau \cdot f)[x_{\sigma(1)}, \ldots, x_{\sigma(N)}]
&= \sum f_{(a_1, \ldots, a_N)} x_{\sigma(\tau(1))}^{a_1} \cdots x_{\sigma(\tau(N))}^{a_N} \\
&= \sum f_{(a_1, \ldots, a_N)} x_{(\sigma\tau)(1)}^{a_1} \cdots x_{(\sigma\tau)(N)}^{a_N}
= (\sigma\tau) \cdot f.
\end{align*}
Comparing with~\eqref{pf.prop.sf.SN-acts.b.1}, we obtain
$(\sigma\tau) \cdot f = \sigma \cdot (\tau \cdot f)$.
\end{proof}

\subsection{The action is by algebra automorphisms}

\begin{proposition}
\label{prop.sf.SN-acts-by-alg-auts}
\lean{AlgebraicCombinatorics.SymmetricPolynomials.permAutomorphism, AlgebraicCombinatorics.SymmetricPolynomials.permAutomorphismHom}
\leantarget
\leanok
The group $S_N$ acts on $\mathcal{P}$ by
$K$-algebra automorphisms. In other words, for each $\sigma \in S_N$, the map
\begin{align*}
\mathcal{P} &\to \mathcal{P}, \\
f &\mapsto \sigma \cdot f
\end{align*}
is a $K$-algebra automorphism of $\mathcal{P}$.
\end{proposition}

\begin{proof}
Fix $\sigma \in S_N$. For any $f, g \in \mathcal{P}$, we have
\begin{equation}
\sigma \cdot (fg) = (\sigma \cdot f) \cdot (\sigma \cdot g)
\label{pf.prop.sf.SN-acts-by-alg-auts.1}
\end{equation}
(by expanding the substitution definition of the action).
Thus, the map $f \mapsto \sigma \cdot f$
respects multiplication. Similarly, this map respects addition, respects
scaling, respects the zero and respects the unity. Hence, this map is a
$K$-algebra morphism from $\mathcal{P}$ to $\mathcal{P}$. Furthermore, this
map is invertible, since its inverse is the map $f \mapsto \sigma^{-1} \cdot f$.
Thus, this map is a $K$-algebra automorphism of $\mathcal{P}$.
\end{proof}

\subsection{Symmetric polynomials form a subalgebra}

\begin{theorem}
\label{thm.sf.S-subalg}
\lean{AlgebraicCombinatorics.SymmetricPolynomials.S_subalgebra}
\leantarget
\leanok
The subset $\mathcal{S}$ is a $K$-subalgebra of $\mathcal{P}$.
\end{theorem}

\begin{proof}
\leanok
We need to show that $\mathcal{S}$ is closed under addition, multiplication and scaling, and that
$\mathcal{S}$ contains the zero and the unity of $\mathcal{P}$. Let us show
that $\mathcal{S}$ is closed under multiplication (since all the other
claims are equally easy): Let $f, g \in \mathcal{S}$. We must show that
$fg \in \mathcal{S}$.

The polynomial $f$ is symmetric (since $f \in \mathcal{S}$); in other words,
$\sigma \cdot f = f$ for each $\sigma \in S_N$. Similarly, $\sigma \cdot g = g$ for
each $\sigma \in S_N$. Now, from~\eqref{pf.prop.sf.SN-acts-by-alg-auts.1}, we
see that
for each $\sigma \in S_N$, we have
$\sigma \cdot (fg) = (\sigma \cdot f) \cdot (\sigma \cdot g) = fg$.
This shows that $fg$ is symmetric, i.e., $fg \in \mathcal{S}$.
The other claims are proved similarly.
\end{proof}

\begin{definition}
\label{def.sf.ring-of-symm}
\lean{AlgebraicCombinatorics.SymmetricPolynomials.S}
\leantarget
\leanok
The $K$-subalgebra $\mathcal{S}$ of $\mathcal{P}$
is called the \emph{ring of symmetric polynomials in} $N$ \emph{variables over} $K$.
\end{definition}

\subsection{Monomials}

\begin{definition}
\label{def.sf.monomial}
\lean{AlgebraicCombinatorics.SymmetricPolynomials.Monomial, AlgebraicCombinatorics.SymmetricPolynomials.Monomial.degree, AlgebraicCombinatorics.SymmetricPolynomials.Monomial.IsSquarefree, AlgebraicCombinatorics.SymmetricPolynomials.Monomial.IsPrimal}
\leantarget
\leanok
\textbf{(a)} A \emph{monomial} is an expression of the
form $x_1^{a_1} x_2^{a_2} \cdots x_N^{a_N}$ with $a_1, a_2, \ldots, a_N \in \mathbb{N}$. \medskip

\textbf{(b)} The \emph{degree} $\deg \mathfrak{m}$
of a monomial $\mathfrak{m} = x_1^{a_1} x_2^{a_2} \cdots x_N^{a_N}$
is defined to be $a_1 + a_2 + \cdots + a_N \in \mathbb{N}$. \medskip

\textbf{(c)} A monomial $\mathfrak{m} = x_1^{a_1} x_2^{a_2} \cdots x_N^{a_N}$
is said to be \emph{squarefree} if $a_1, a_2, \ldots, a_N \in \{0,1\}$.
(This is saying that no square or higher power of
an indeterminate appears in $\mathfrak{m}$; thus the name ``squarefree''.) \medskip

\textbf{(d)} A monomial $\mathfrak{m} = x_1^{a_1} x_2^{a_2} \cdots x_N^{a_N}$
is said to be \emph{primal} if there is at most one
$i \in [N]$ satisfying $a_i > 0$. (This is saying that the
monomial $\mathfrak{m}$ contains no two distinct indeterminates. Thus, a
primal monomial is just $1$ or a power of an indeterminate.)
\end{definition}

\subsection{Elementary, complete homogeneous, and power sum symmetric polynomials}

\begin{definition}
\label{def.sf.ehp}
\lean{AlgebraicCombinatorics.SymmetricPolynomials.e, AlgebraicCombinatorics.SymmetricPolynomials.h, AlgebraicCombinatorics.SymmetricPolynomials.p}
\leantarget
\leanok
\textbf{(a)} For each $n \in \mathbb{Z}$, define a symmetric
polynomial $e_n \in \mathcal{S}$ by
\[
e_n = \sum_{\substack{(i_1, i_2, \ldots, i_n) \in [N]^n; \\ i_1 < i_2 < \cdots < i_n}}
x_{i_1} x_{i_2} \cdots x_{i_n}
= (\text{sum of all squarefree monomials of degree } n).
\]
This $e_n$ is called the $n$-th \emph{elementary symmetric polynomial} in
$x_1, x_2, \ldots, x_N$. \medskip

\textbf{(b)} For each $n \in \mathbb{Z}$, define a symmetric polynomial
$h_n \in \mathcal{S}$ by
\[
h_n = \sum_{\substack{(i_1, i_2, \ldots, i_n) \in [N]^n; \\ i_1 \leq i_2 \leq \cdots \leq i_n}}
x_{i_1} x_{i_2} \cdots x_{i_n}
= (\text{sum of all monomials of degree } n).
\]
This $h_n$ is called the $n$-th \emph{complete homogeneous symmetric polynomial}
in $x_1, x_2, \ldots, x_N$. \medskip

\textbf{(c)} For each $n \in \mathbb{Z}$, define a symmetric polynomial
$p_n \in \mathcal{S}$ by
\begin{align*}
p_n &=
\begin{cases}
x_1^n + x_2^n + \cdots + x_N^n, & \text{if } n > 0; \\
1, & \text{if } n = 0; \\
0, & \text{if } n < 0
\end{cases}
\\
&= (\text{sum of all primal monomials of degree } n).
\end{align*}
This $p_n$ is called the $n$-th \emph{power sum} in $x_1, x_2, \ldots, x_N$.
\end{definition}

\subsection{Elementary symmetric polynomials vanish for large degree}

\begin{proposition}
\label{prop.sf.en=0}
\lean{AlgebraicCombinatorics.SymmetricPolynomials.e_eq_zero_of_gt}
\leantarget
\leanok
For each integer $n > N$, we have $e_n = 0$.
\end{proposition}

\begin{proof}
\leanok
Let $n > N$ be an integer. Then, the set $[N]$ has no $n$
distinct elements. Thus, there exists no $n$-tuple $(i_1, i_2, \ldots, i_n)
\in [N]^n$ satisfying $i_1 < i_2 < \cdots < i_n$.
Now, the definition of $e_n$ yields
\begin{align*}
e_n &= \sum_{\substack{(i_1, \ldots, i_n) \in [N]^n; \\ i_1 < \cdots < i_n}}
x_{i_1} \cdots x_{i_n}
= (\text{empty sum}) = 0.
\end{align*}
\end{proof}

\subsection{Generating function identities}

\begin{proposition}
\label{prop.sf.e-h-FPS}
\lean{AlgebraicCombinatorics.SymmetricPolynomials.esymm_genfunc, AlgebraicCombinatorics.SymmetricPolynomials.esymm_genfunc_two_var, AlgebraicCombinatorics.SymmetricPolynomials.hsymm_genfunc}
\leantarget
\leanok
\textbf{(a)} In the polynomial ring $\mathcal{P}[t]$, we have
\[
\prod_{i=1}^{N} (1 - t x_i) = \sum_{n \in \mathbb{N}} (-1)^n t^n e_n.
\]

\textbf{(b)} In the polynomial ring $\mathcal{P}[u,v]$, we have
\[
\prod_{i=1}^{N} (u - v x_i) = \sum_{n=0}^{N} (-1)^n u^{N-n} v^n e_n.
\]

\textbf{(c)} In the FPS ring $\mathcal{P}[[t]]$, we have
\[
\prod_{i=1}^{N} \frac{1}{1 - t x_i} = \sum_{n \in \mathbb{N}} t^n h_n.
\]
\end{proposition}

\begin{proof}
\textbf{(a)} For each $i \in \{1, 2, \ldots, N\}$, we have
$1 + t x_i = \sum_{a \in \{0,1\}} (t x_i)^a$.
Multiplying these equalities over all $i \in \{1, 2, \ldots, N\}$, we obtain
\begin{align*}
\prod_{i=1}^{N} (1 + t x_i)
&= \sum_{(a_1, \ldots, a_N) \in \{0,1\}^N} t^{a_1 + \cdots + a_N} x_1^{a_1} \cdots x_N^{a_N} \\
&= \sum_{\substack{\mathfrak{m} \text{ squarefree}}} t^{\deg \mathfrak{m}} \mathfrak{m}
= \sum_{n \in \mathbb{N}} t^n e_n.
\end{align*}
Substituting $-t$ for $t$ gives the stated identity. \medskip

\textbf{(b)} This is similar to part \textbf{(a)}, using two variables $u$ and $v$.  \medskip

\textbf{(c)} For each $i \in \{1, 2, \ldots, N\}$, we have
$\frac{1}{1 - t x_i} = \sum_{a \in \mathbb{N}} (t x_i)^a$.
Multiplying these equalities over all $i \in \{1, 2, \ldots, N\}$, we obtain
\begin{align*}
\prod_{i=1}^{N} \frac{1}{1 - t x_i}
&= \sum_{(a_1, \ldots, a_N) \in \mathbb{N}^N} t^{a_1 + \cdots + a_N} x_1^{a_1} \cdots x_N^{a_N} \\
&= \sum_{\mathfrak{m} \text{ monomial}} t^{\deg \mathfrak{m}} \mathfrak{m}
= \sum_{n \in \mathbb{N}} t^n h_n.
\end{align*}
\end{proof}

\subsubsection{Helpers for the generating function proofs}

\begin{lemma}
\label{lem.sf.geom-series}
\lean{AlgebraicCombinatorics.SymmetricPolynomials.geom_series_mul_one_sub, AlgebraicCombinatorics.SymmetricPolynomials.prod_geom_series_mul_prod_one_sub}
\leanhelper
\leanok
For each $i$, the geometric series identity
$(1 - t x_i) \cdot \sum_{k \geq 0} (t x_i)^k = 1$
holds in the power series ring $\mathcal{P}[[t]]$.
Taking the product over all $i \in [N]$, we obtain the generating function identity
$E(t) \cdot H(t) = 1$, where
$E(t) = \prod_{i=1}^N (1 - t x_i)$ and
$H(t) = \prod_{i=1}^N \sum_{k \geq 0} (t x_i)^k$.
\end{lemma}

\begin{proof}
The identity $(1 - t x_i) \cdot \sum_{k \geq 0} (t x_i)^k = 1$ is proved by extracting
coefficients: the coefficient of $t^n$ on the left is $x_i^n - x_i \cdot x_i^{n-1} = 0$ for $n > 0$
and $1$ for $n = 0$. The product identity follows from
$\prod_i (1 - t x_i) \cdot \prod_i \sum_k (t x_i)^k = \prod_i 1 = 1$.
\end{proof}

\subsection{Newton--Girard formulas}

\begin{theorem}[Newton--Girard formulas]
\label{thm.sf.NG}
\lean{AlgebraicCombinatorics.SymmetricPolynomials.newtonGirard_eh, AlgebraicCombinatorics.SymmetricPolynomials.newtonGirard_hp, AlgebraicCombinatorics.SymmetricPolynomials.newtonGirard_esymm}
\leantarget
\leanok
For any positive integer $n$, we have
\begin{align}
\sum_{j=0}^{n} (-1)^j e_j h_{n-j} &= 0; \label{eq.thm.sf.NG.eh} \\
\sum_{j=1}^{n} (-1)^{j-1} e_{n-j} p_j &= n e_n; \label{eq.thm.sf.NG.ep} \\
\sum_{j=1}^{n} h_{n-j} p_j &= n h_n. \label{eq.thm.sf.NG.hp}
\end{align}
\end{theorem}

\begin{proof}
We prove formulas~\eqref{eq.thm.sf.NG.eh} and~\eqref{eq.thm.sf.NG.hp}; the formula~\eqref{eq.thm.sf.NG.ep}
can be derived by differentiating the generating function
$\prod_{i=1}^N (1 - t x_i) = \sum_n (-1)^n t^n e_n$
with respect to $t$, and comparing coefficients (see the source for details).

\textbf{Proof of~\eqref{eq.thm.sf.NG.eh}:}
In the FPS ring $\mathcal{P}[[t]]$, by Proposition~\ref{prop.sf.e-h-FPS}~(a) and~(c) we have
\[
\prod_{i=1}^{N} (1 - t x_i) = \sum_{n \in \mathbb{N}} (-1)^n t^n e_n
\]
and
\[
\prod_{i=1}^{N} \frac{1}{1 - t x_i} = \sum_{n \in \mathbb{N}} t^n h_n.
\]
Multiplying these two equalities, the left-hand side becomes
\[
\left(\prod_{i=1}^{N} (1 - t x_i)\right) \left(\prod_{i=1}^{N} \frac{1}{1 - t x_i}\right)
= \prod_{i=1}^{N} 1 = 1.
\]
The right-hand side product is
$\sum_{n \in \mathbb{N}} \left(\sum_{j=0}^n (-1)^j e_j h_{n-j}\right) t^n$.
Comparing coefficients of $t^n$ for $n > 0$, we obtain~\eqref{eq.thm.sf.NG.eh}.

\textbf{Proof of~\eqref{eq.thm.sf.NG.hp}:}
The proof is by a counting argument on monomials.
Each monomial in $h_n$ corresponds to a multiset $s$ of size $n$ from $[N]$.
On the left-hand side, each such monomial appears exactly $n$ times:
once for each way to pick an element $i$ in the support of $s$
and a positive integer $j \leq \text{count}(i, s)$,
contributing to the term $h_{n-j} p_j$.
For each $s$, the total count is $\sum_{i=1}^N \text{count}(i, s) = n$, matching the right-hand side.
\end{proof}

\subsection{Fundamental Theorem of Symmetric Polynomials}

\begin{theorem}[Fundamental Theorem of Symmetric Polynomials]
\label{thm.sf.ftsf}
\lean{AlgebraicCombinatorics.SymmetricPolynomials.esymmAlgEquiv', AlgebraicCombinatorics.SymmetricPolynomials.esymm_algebraicIndependent, AlgebraicCombinatorics.SymmetricPolynomials.esymm_generates_symmetric, AlgebraicCombinatorics.SymmetricPolynomials.hsymm_algebraicIndependent, AlgebraicCombinatorics.SymmetricPolynomials.hsymm_generates_symmetric, AlgebraicCombinatorics.SymmetricPolynomials.psum_algebraicIndependent, AlgebraicCombinatorics.SymmetricPolynomials.psum_generates_symmetric}
\leantarget
\leanok
\textbf{(a)} The elementary symmetric polynomials $e_1, e_2, \ldots, e_N$
are algebraically independent (over $K$) and generate the $K$-algebra $\mathcal{S}$.

In other words, each $f \in \mathcal{S}$ can be uniquely written as a polynomial
in $e_1, e_2, \ldots, e_N$.

In yet other words, the map
\begin{align*}
K[y_1, y_2, \ldots, y_N] &\to \mathcal{S}, \\
g &\mapsto g[e_1, e_2, \ldots, e_N]
\end{align*}
is a $K$-algebra isomorphism. \medskip

\textbf{(b)} The complete homogeneous symmetric polynomials $h_1, h_2, \ldots, h_N$
are algebraically independent (over $K$) and generate the $K$-algebra $\mathcal{S}$.

In other words, each $f \in \mathcal{S}$ can be uniquely written as a polynomial
in $h_1, h_2, \ldots, h_N$.

In yet other words, the map
\begin{align*}
K[y_1, y_2, \ldots, y_N] &\to \mathcal{S}, \\
g &\mapsto g[h_1, h_2, \ldots, h_N]
\end{align*}
is a $K$-algebra isomorphism. \medskip

\textbf{(c)} Now assume that $K$ is a commutative $\mathbb{Q}$-algebra (e.g.,
a field of characteristic $0$). Then, the power sums $p_1, p_2, \ldots, p_N$
are algebraically independent (over $K$) and generate the $K$-algebra $\mathcal{S}$.

In other words, each $f \in \mathcal{S}$ can be uniquely written as a polynomial
in $p_1, p_2, \ldots, p_N$.

In yet other words, the map
\begin{align*}
K[y_1, y_2, \ldots, y_N] &\to \mathcal{S}, \\
g &\mapsto g[p_1, p_2, \ldots, p_N]
\end{align*}
is a $K$-algebra isomorphism.
\end{theorem}

\begin{proof}
Various proofs can be found in \cite[Theorem 171]{Benede25}, \cite[proof of Theorem 1]{BluCos16},
\cite[Theorem 1.2.1]{Dumas08}, \cite[Theorem 9.6.6]{Goodman}, \cite[\S 1.1]{Smith95}.

\textbf{Part (a):}
Part (a) is proved by constructing the evaluation map
$K[y_1, \ldots, y_N] \to \mathcal{S}$ sending $y_k \mapsto e_k$
and showing it is a $K$-algebra isomorphism.
Algebraic independence follows from injectivity;
generation follows from surjectivity.

\textbf{Part (b):}
The proof uses the Newton--Girard relations to show that each $e_k$ (for $k \leq N$)
can be expressed as a polynomial in $h_1, \ldots, h_N$, by strong induction on $k$.
The key identity is: from~\eqref{eq.thm.sf.NG.eh},
$(-1)^{k+1} e_{k+1} = -\sum_{j=0}^{k} (-1)^j e_j h_{k+1-j}$,
where by the induction hypothesis $e_0, \ldots, e_k$ are already in the range of the
evaluation map at $h_1, \ldots, h_N$. The algebraic independence of the $h_i$ then
follows by a transcendence degree argument: the composition
$K[y_1, \ldots, y_N] \to \mathcal{S} \xrightarrow{\sim} K[y_1, \ldots, y_N]$
(via the isomorphism from part (a)) is a surjective endomorphism, hence bijective
since polynomial rings over integral domains have finite transcendence degree.

Note: The Lean proof of part (b) requires the additional hypothesis that $K$ is an integral domain.

\textbf{Part (c):}
The argument is analogous to part (b). Newton's identity~\eqref{eq.thm.sf.NG.ep}
gives $n \cdot e_n = \sum_{j=1}^n (-1)^{j-1} e_{n-j} p_j$. Since $K$ is a
$\mathbb{Q}$-algebra, $n$ is invertible in $K$, so we can solve for $e_n$ as a
polynomial in $e_0, \ldots, e_{n-1}$ and $p_1, \ldots, p_n$. By strong induction,
each $e_k$ lies in the range of the evaluation map at $p_1, \ldots, p_N$.
Algebraic independence again follows by the transcendence degree argument.

Note: The Lean formalization of part (c) contains sorry in some helper lemmas
(weight constraint and triangular structure arguments).
\end{proof}

\subsection{Simple transpositions suffice}

\begin{lemma}
\label{lem.sf.simples-enough}
\lean{AlgebraicCombinatorics.SymmetricPolynomials.isSymm_iff_simpleTranspositions}
\leantarget
\leanok
For each $i \in [N-1]$, let $s_i \in S_N$ denote the simple transposition
that swaps $i$ and $i+1$.

Let $f \in \mathcal{P}$. Assume that
\begin{equation}
s_k \cdot f = f \quad \text{for each } k \in [N-1].
\label{eq.lem.sf.simples-enough.ass}
\end{equation}
Then, the polynomial $f$ is symmetric.
\end{lemma}

\begin{proof}
The proof proceeds by showing that invariance under simple transpositions
implies invariance under all transpositions of the form $\text{swap}(i, j)$
for any $i, j \in [N]$.

\textbf{Step 1:} Invariance under adjacent transpositions
implies invariance under arbitrary transpositions.
We prove by induction on $|j - i|$ that $\text{swap}(i, j) \cdot f = f$.
The base case $|j - i| = 1$ is the hypothesis~\eqref{eq.lem.sf.simples-enough.ass}.
For the inductive step, use the conjugation identity:
$\text{swap}(i, j) = \text{swap}(i, i+1) \cdot \text{swap}(i+1, j) \cdot \text{swap}(i, i+1)$.

\textbf{Step 2:} Every permutation $\sigma \in S_N$ can be written as a product
of transpositions (this is a standard fact about the symmetric group).
Since each transposition fixes $f$, so does $\sigma$.
\end{proof}

\subsection{Helper lemmas from the Lean formalization}

The following results appear in the Lean formalization but not in the TeX source.
They provide API for the definitions above.

\begin{lemma}
\label{lem.sf.monomial-degree-le-squarefree}
\lean{AlgebraicCombinatorics.SymmetricPolynomials.Monomial.degree_le_of_isSquarefree}
\leanhelper
\leanok
A squarefree monomial has degree at most $N$.
\end{lemma}

\begin{proof}
\leanok
If $\mathfrak{m} = x_1^{a_1} \cdots x_N^{a_N}$ is squarefree, then
each $a_i \leq 1$, so $\deg \mathfrak{m} = a_1 + \cdots + a_N \leq N$.
\end{proof}

\begin{lemma}
\label{lem.sf.sum-count-eq-card}
\lean{AlgebraicCombinatorics.SymmetricPolynomials.sum_count_eq_card}
\leanhelper
\leanok
For any multiset $s$ of elements from $[N]$ of size $m$,
the sum of counts $\sum_{i=1}^N \text{count}(i, s)$ equals $m$.
\end{lemma}

\begin{proof}
This is the standard identity that the sum of all multiplicities in
a multiset equals its cardinality. The proof splits the sum over
elements that appear in $s$ and elements that do not, the latter contributing zero.
\end{proof}

\begin{lemma}
\label{lem.sf.e-isSymm}
\lean{AlgebraicCombinatorics.SymmetricPolynomials.e_isSymm, AlgebraicCombinatorics.SymmetricPolynomials.h_isSymm, AlgebraicCombinatorics.SymmetricPolynomials.p_isSymm}
\leanhelper
\leanok
The elementary symmetric polynomial $e_n$, the complete homogeneous symmetric
polynomial $h_n$, and the power sum $p_n$ are all symmetric (i.e., belong to
$\mathcal{S}$) for every $n$.
\end{lemma}

\begin{proof}
\leanok
For $e_n$: permuting the variables permutes the strictly increasing
$n$-tuples $(i_1 < \cdots < i_n)$, so the sum is unchanged.
For $h_n$ and $p_n$: the argument is analogous, using the fact that
permuting variables permutes the non-decreasing tuples (for $h_n$)
or the individual terms $x_i^n$ (for $p_n$).
\end{proof}

\begin{lemma}
\label{lem.sf.monomial-isPrimal-iff}
\lean{AlgebraicCombinatorics.SymmetricPolynomials.Monomial.isPrimal_iff}
\leanhelper
\leanok
A monomial $\mathfrak{m}$ is primal if and only if $\mathfrak{m} = 1$
or $\mathfrak{m} = x_i^k$ for some $i \in [N]$ and some $k > 0$.
\end{lemma}

\begin{proof}
\leanok
If $\mathfrak{m}$ is primal, its support has at most one element.
If the support is empty, then $\mathfrak{m} = 1$.
If the support is $\{i\}$, then $\mathfrak{m} = x_i^{m_i}$ for some $m_i > 0$.
The converse is immediate.
\end{proof}

\subsection{Monomial API}

The following lemmas develop the API for monomials, connecting the exponent vector
representation with the polynomial representation.

\begin{lemma}
\label{lem.sf.monomial-toPoly}
\lean{AlgebraicCombinatorics.SymmetricPolynomials.Monomial.toPoly}
\leanhelper
\leanok
A monomial $\mathfrak{m}$ with exponent vector $(a_1, \ldots, a_N)$ corresponds to the
polynomial $x_1^{a_1} \cdots x_N^{a_N}$.
The polynomial corresponding to the zero exponent vector is $1$,
the polynomial corresponding to the $i$-th standard basis vector is $x_i$,
and the polynomial corresponding to a sum of exponent vectors $m_1 + m_2$ is
the product of the corresponding polynomials.
\end{lemma}

\begin{proof}
Direct from the definition of monomials in a multivariable polynomial ring.
\end{proof}

\begin{lemma}
\label{lem.sf.monomial-ofFinset}
\lean{AlgebraicCombinatorics.SymmetricPolynomials.Monomial.ofFinset, AlgebraicCombinatorics.SymmetricPolynomials.Monomial.toPoly_ofFinset, AlgebraicCombinatorics.SymmetricPolynomials.Monomial.isSquarefree_ofFinset, AlgebraicCombinatorics.SymmetricPolynomials.Monomial.degree_ofFinset, AlgebraicCombinatorics.SymmetricPolynomials.Monomial.support_ofFinset}
\leanhelper
\leanok
For a subset $S \subseteq [N]$, the monomial $\mathfrak{m}_S = \prod_{i \in S} x_i$
is squarefree, has degree $|S|$, and has support equal to $S$.
\end{lemma}

\begin{proof}
\leanok
The monomial $\mathfrak{m}_S$ is defined as $\sum_{i \in S} e_i$ where $e_i$ is the
$i$-th standard basis vector. Each entry is $0$ or $1$ (squarefree),
the sum of entries equals $|S|$ (degree), and the support is exactly $S$.
All properties are proved by induction on $|S|$.
\end{proof}

\begin{lemma}
\label{lem.sf.monomial-degree-add}
\lean{AlgebraicCombinatorics.SymmetricPolynomials.Monomial.degree_add}
\leanhelper
\leanok
The degree of the product of two monomials is the sum of their degrees:
$\deg(\mathfrak{m}_1 \cdot \mathfrak{m}_2) = \deg \mathfrak{m}_1 + \deg \mathfrak{m}_2$.
\end{lemma}

\begin{proof}
\leanok
The degree is defined as the sum of all exponents. Since the exponent vector
of $\mathfrak{m}_1 \cdot \mathfrak{m}_2$ is the pointwise sum of the
exponent vectors, the degree of the product equals the sum of the degrees.
\end{proof}

\subsection{Basic values of $e_n$, $h_n$, $p_n$}

\begin{lemma}
\label{lem.sf.basic-values}
\lean{AlgebraicCombinatorics.SymmetricPolynomials.e_zero, AlgebraicCombinatorics.SymmetricPolynomials.e_one, AlgebraicCombinatorics.SymmetricPolynomials.h_zero, AlgebraicCombinatorics.SymmetricPolynomials.h_one, AlgebraicCombinatorics.SymmetricPolynomials.p_zero, AlgebraicCombinatorics.SymmetricPolynomials.p_one}
\leanhelper
\leanok
We have the following basic values:
\begin{itemize}
\item $e_0 = h_0 = 1$;
\item $e_1 = h_1 = p_1 = \sum_{i=1}^N x_i$;
\item $p_0 = N$ (the number of variables; note this differs from the textbook convention $p_0 = 1$).
\end{itemize}
\end{lemma}

\begin{proof}
These follow directly from the definitions. For $e_0$ and $h_0$, the only
$0$-tuple is the empty tuple, giving an empty product equal to $1$. For $e_1$
and $h_1$, the $1$-tuples are the singletons $(i)$ for $i \in [N]$, giving
$\sum_{i=1}^N x_i$. For $p_1$, this is $\sum_{i=1}^N x_i^1 = \sum_{i=1}^N x_i$.
For $p_0$, the Lean definition follows Mathlib's convention, which gives
$p_0 = \sum_{i=1}^N x_i^0 = N$.
\end{proof}

\subsection{Alternative characterizations of $e_n$, $h_n$, $p_n$}

\begin{lemma}
\label{lem.sf.ehp-alt-characterizations}
\lean{AlgebraicCombinatorics.SymmetricPolynomials.e_eq_sum_prod_subsets, AlgebraicCombinatorics.SymmetricPolynomials.p_eq_sum_pow}
\leanhelper
\leanok
The defining formulas for $e_n$ and $p_n$ are:
\[
e_n = \sum_{\substack{S \subseteq [N] \\ |S| = n}} \prod_{i \in S} x_i,
\qquad
p_n = \sum_{i=1}^N x_i^n.
\]
\end{lemma}

\begin{proof}
Both follow immediately from the definitions:
$e_n$ is the sum over strictly increasing $n$-tuples, which bijects with
$n$-element subsets $S \subseteq [N]$; and $p_n = \sum_{i=1}^N x_i^n$
is the definition for $n > 0$.
\end{proof}

\subsection{Integer-indexed versions}

\begin{lemma}
\label{lem.sf.integer-indexed}
\lean{AlgebraicCombinatorics.SymmetricPolynomials.eZ, AlgebraicCombinatorics.SymmetricPolynomials.hZ, AlgebraicCombinatorics.SymmetricPolynomials.pZ}
\leanhelper
\leanok
The integer-indexed versions $e_n$, $h_n$, $p_n$ for $n \in \mathbb{Z}$ satisfy:
\begin{itemize}
\item $e_n = 0$ for $n < 0$, and $e_n$ agrees with Definition~\ref{def.sf.ehp}~\textbf{(a)} for $n \geq 0$;
\item $h_n = 0$ for $n < 0$, and $h_n$ agrees with Definition~\ref{def.sf.ehp}~\textbf{(b)} for $n \geq 0$;
\item $p_n = 0$ for $n < 0$, $p_0 = 1$, and $p_n = x_1^n + x_2^n + \cdots + x_N^n$ for $n > 0$.
\end{itemize}
All are symmetric for every $n \in \mathbb{Z}$.
\end{lemma}

\begin{proof}
By case analysis on the sign of $n$.
For non-negative $n$, symmetry follows from Lemma~\ref{lem.sf.e-isSymm}.
For negative $n$, the polynomial is $0$ which is trivially symmetric.
\end{proof}

\subsection{Adding a variable recurrence}

\begin{lemma}
\label{lem.sf.esymm-add-var}
\lean{AlgebraicCombinatorics.SymmetricPolynomials.esymm_succ_add_var', AlgebraicCombinatorics.SymmetricPolynomials.esymm_succ_add_var}
\leanhelper
\leanok
For each $n \geq 0$, the elementary symmetric polynomial satisfies the recurrence
\[
e_{n+1}(x_1, \ldots, x_N, y) = e_{n+1}(x_1, \ldots, x_N) + y \cdot e_n(x_1, \ldots, x_N).
\]
This follows by partitioning the $n$-element subsets of $\{1, \ldots, N+1\}$ into those
containing $N+1$ and those not.
\end{lemma}

\begin{proof}
Partition $\binom{[N+1]}{n+1}$ into subsets not containing $N+1$ and subsets containing $N+1$.
The first group yields $e_{n+1}(x_1, \ldots, x_N)$ (via a bijection with $\binom{[N]}{n+1}$).
The second group yields $y \cdot e_n(x_1, \ldots, x_N)$ (via a bijection with $\binom{[N]}{n}$,
factoring out the variable $x_{N+1} = y$). The proof establishes
these bijections explicitly.
\end{proof}

\subsection{Antisymmetric polynomials}

\begin{definition}
\label{def.sf.antisymm}
\lean{AlgebraicCombinatorics.SymmetricPolynomials.IsAntisymm}
\leanhelper
\leanok
A polynomial $f \in \mathcal{P}$ is \emph{antisymmetric} if
\[
\sigma \cdot f = \operatorname{sign}(\sigma) \cdot f \quad \text{for all } \sigma \in S_N.
\]
\end{definition}

\begin{lemma}
\label{lem.sf.antisymm-submodule}
\lean{AlgebraicCombinatorics.SymmetricPolynomials.isAntisymm_zero, AlgebraicCombinatorics.SymmetricPolynomials.isAntisymm_add, AlgebraicCombinatorics.SymmetricPolynomials.isAntisymm_smul, AlgebraicCombinatorics.SymmetricPolynomials.isAntisymm_neg}
\leanhelper
\leanok
The antisymmetric polynomials form a $K$-submodule of $\mathcal{P}$: the zero polynomial
is antisymmetric, and the sum, scalar multiple, and negation of antisymmetric polynomials
are antisymmetric.
\end{lemma}

\begin{proof}
\leanok
For each property, apply a permutation $\sigma$ and use the linearity of the action.
For instance, $\sigma \cdot (f + g) = \sigma \cdot f + \sigma \cdot g
= \operatorname{sign}(\sigma)(f + g)$.
\end{proof}

\begin{lemma}
\label{lem.sf.antisymm-sq-symmetric}
\lean{AlgebraicCombinatorics.SymmetricPolynomials.isSymm_sq_of_isAntisymm}
\leanhelper
\leanok
The square of an antisymmetric polynomial is symmetric.
\end{lemma}

\begin{proof}
\leanok
If $f$ is antisymmetric, then
$\sigma \cdot (f^2) = (\sigma \cdot f)^2 = (\operatorname{sign}(\sigma) \cdot f)^2
= \operatorname{sign}(\sigma)^2 \cdot f^2 = f^2$,
using $\operatorname{sign}(\sigma)^2 = 1$.
\end{proof}

\subsection{The sum of variables is symmetric}

\begin{lemma}
\label{lem.sf.sum-X-symmetric}
\lean{AlgebraicCombinatorics.SymmetricPolynomials.isSymm_sum_X}
\leanhelper
\leanok
The polynomial $\sum_{i=1}^N x_i$ is symmetric.
\end{lemma}

\begin{proof}
\leanok
This equals $e_1$, which is symmetric by Lemma~\ref{lem.sf.e-isSymm}.
\end{proof}
